\section{Resultados y discusión}
\subsection{Alcance de las simulaciones}
Se ejecutaron en MATLAB R2026b las cuatro configuraciones nominales
(predicción y estado actual; ubicaciones E1 y E2) con
$T_s=\SI{2}{\milli\second}$ y referencia de \SI{50}{\milli\volt}
aplicada en $t=\SI{20}{\milli\second}$.
La salida $y[k]$ es el único estado medido que se realimentaría físicamente.
Los valores de esta sección son \emph{exclusivamente simulaciones}, no
mediciones en PSoC. Las gráficas E2 permanecen pendientes de discusión.

\begin{table}[t]
\centering
\caption{Estado de los casos del Laboratorio 7.}
\label{tab:ensayos}
\small
\begin{tabular}{@{}llll@{}}
\toprule
Observador & Polos & MATLAB & PSoC \\
\midrule
Predicción & E1 & Analizado & Pendiente \\
Predicción & E2 & Ejecutado & Pendiente \\
Actual & E1 & Analizado & Pendiente \\
Actual & E2 & Ejecutado & Pendiente \\
\bottomrule
\end{tabular}
\end{table}

\subsection{Comparación nominal entre observadores E1}
Para ambos controladores se utilizaron los polos de lazo cerrado
$0.8\pm j0.25$, los polos de observación E1 $0.4\pm j0.25$
y las mismas condiciones iniciales: $x[0]=[0\quad0]^\top$,
$\hat{x}[0]=[0.02\quad-0.01]^\top$. Las métricas de la
Tabla~\ref{tab:res_e1} se calcularon reproduciendo numéricamente el
algoritmo de las simulaciones MATLAB, y se contrastaron con las
gráficas aportadas del ensayo computacional.

\begin{table}[t]
\centering
\caption{Comparación E1: simulación nominal, no ensayo físico.}
\label{tab:res_e1}
\small
\begin{tabular}{@{}lrr@{}}
\toprule
Magnitud & Predictor & Actual \\
\midrule
Máximo de salida (\si{\milli\volt}) & 57.646 & 57.631 \\
Sobrepasamiento (\%) & 15.293 & 15.263 \\
Pico desde el escalón (\si{\milli\second}) & 22 & 22 \\
Establecimiento $\pm2\%$ (\si{\milli\second}) & 46 & 48 \\
$u[0]$ (\si{\milli\volt}) & 25.57 & 55.90 \\
$\max u$ (\si{\milli\volt}) & 97.98 & 97.88 \\
$\min u$ (\si{\milli\volt}) & -8.72 & -16.25 \\
$u_\infty$ (\si{\milli\volt}) & 49.09 & 49.09 \\
\bottomrule
\end{tabular}
\end{table}

El predictor presentó un máximo de \SI{8.20}{\milli\volt}
antes de aplicar la referencia, frente a \SI{15.22}{\milli\volt}
del estimador actual. Esta diferencia surge del error de inicialización:
en el estimador actual, la medición de la muestra $k$ modifica la
estimación antes de calcular $u[k]$; el predictor utiliza la innovación
para construir el estado estimado de la siguiente muestra.
Aunque el estimador actual utiliza información más reciente, no
proporciona por sí mismo un menor tiempo de establecimiento: en este
caso, ambos comportamientos son muy semejantes.

La norma de error inicial reportada por el estimador actual está
calculada \emph{después} de corregir su predicción con $y[0]$,
mientras que la del predictor es previa a esa corrección. Por ello
no debe compararse directamente la altura inicial de las curvas
como si representaran el mismo instante de cálculo.
En condiciones ideales, ambas arquitecturas presentan convergencia
del error; resta estudiar ruido, cuantización y retardos reales.

\subsection{Ensayos experimentales pendientes}
Para cada diseño deberán medirse los tiempos de establecimiento,
los sobrepasamientos, el error estacionario, el máximo y el RMS del
esfuerzo, y el efecto del ruido. La comparación MATLAB--banco requiere
datos exportados del Rigol DS1202Z-E y calibraciones ADC/DAC actuales.
No se atribuyen las métricas nominales al hardware físico.
